
\sgaStableTarget{sga6.e10.sec.1.2}\subsection*{1.2.}
Now let @@M0@@ be an arbitrary scheme, and consider the group @@M1@@ of classes of locally free @@M2@@-modules (henceforth understood to be of finite type). It is an augmented pre-@@M3@@-algebra over the ring @@M4@@ (V 3.9.1). The map @@M5@@ allows us henceforth to identify @@M6@@ with a subring of @@M7@@. If @@M8@@ is a locally free @@M9@@-module, its class in @@M11@@ will be denoted by @@M10@@. Define the operators @@M12@@ (@@M13@@) by
\sgaStableTarget{sga6.e10.eq.1-2-1}@@M14@@
and by
\sgaStableTarget{sga6.e10.eq.1-2-2}@@M15@@
(V 2.2). Recall that in the @@M16@@-filtration @@M17@@ of @@M18@@, one has @@M19@@ for @@M20@@, and that for @@M21@@, @@M22@@ is the ideal generated by the products
\sgaStableTarget{sga6.e10.eq.1-2-3}@@M23@@
with @@M24@@; @@M25@@ and @@M26@@ (@@M27@@ being the augmentation homomorphism, cf. Exposé V 3.10). It was proved (V 3.10) that @@M28@@ is even the @@M29@@-submodule (hence the abelian subgroup, when @@M30@@ is connected) generated by the products (1.2.3), where one may moreover suppose that the elements @@M31@@ (@@M32@@) have the form
\sgaStableTarget{sga6.e10.eq.1-2-4}@@M33@@
with @@M34@@ a locally free @@M35@@-module of rank @@M36@@.

